% !TEX root = task-3-5-4.tex
% Editable task bank. Append new BankTask blocks at the end of the document.
\documentclass[12pt,a4paper]{article}
\usepackage{physics-textbook}
\usepackage{graphicx,needspace}
\usetikzlibrary{decorations.pathmorphing}
\newif\ifShowSolutions
\ifdefined\TutorEdition
  \ShowSolutionstrue
\else
  \ShowSolutionsfalse % Change to \ShowSolutionstrue to display solutions.
\fi
\newcounter{banktask}
\newsavebox{\BankTaskBox}
\newcommand{\BankTaskContent}[3]{%
  \noindent\textbf{\thebanktask.}\quad #2\par
  \ifShowSolutions
    \Needspace{50mm}% Keep the solution heading with its opening text or diagram.
    #3\par
    \ifdefined\TutorEdition
      \ifcsname TutorSource#1\endcsname
        \nopagebreak\smallskip{\small\csname TutorSource#1\endcsname\par}%
      \fi
    \fi
  \fi
}
\newcommand{\BankTask}[3]{%
  \refstepcounter{banktask}\par\medskip
  \sbox{\BankTaskBox}{\begin{minipage}{\linewidth}%
    \BankTaskContent{#1}{#2}{#3}\end{minipage}}%
  \ifdim\dimexpr\ht\BankTaskBox+\dp\BankTaskBox\relax>\dimexpr\textheight-20mm\relax
    \Needspace{35mm}\BankTaskContent{#1}{#2}{#3}%
  \else
    \Needspace{\dimexpr\ht\BankTaskBox+\dp\BankTaskBox\relax}%
    \noindent\usebox{\BankTaskBox}\par
  \fi
  \medskip
}
\begin{document}
\BookSubsection{3.5.4}{Потенциал проводника}


% Source sdamgia-6359: https://phys-ege.sdamgia.ru/problem?id=6359
\BankTask{sdamgia-6359}{
Сплошной металлический шар радиусом $R$, находящийся в вакууме, имеет заряд $Q$. Установите соответствие между физическими величинами и формулами, по которым их можно рассчитать.\par

\begin{center}\begin{minipage}[t]{.47\linewidth}\textbf{Физические величины}\par\smallskip А) модуль напряженности электрического поля на расстоянии $2R$ от центра шара\par
Б) потенциал поверхности шара\end{minipage}\hfill\begin{minipage}[t]{.47\linewidth}\textbf{Формулы}\par\smallskip $1$) $\frac{ 1}{ 4 \pi \varepsilon_{0} } \cdot \frac{ Q}{ 2R }$ \par
$2$) $0$\par
$3$) $\frac{ 1}{ 4 \pi \varepsilon_{0} } \cdot \frac{ Q}{ R }$ \par
$4$) $\frac{ 1}{ 4 \pi \varepsilon_{0} } \cdot \frac{ Q}{ 4R ^2 }$\end{minipage}\end{center}
\par
\begin{center}\small\setlength{\tabcolsep}{3pt}\begin{tabular}{|p{0.45\linewidth}|p{0.45\linewidth}|}\hline
$A$ & Б\\\hline
 & \\\hline
\end{tabular}\end{center}\par
}{
\smallskip\noindent\textbf{Решение}\par
А) Модуль напряженности сферически симметричного распределения заряда вне этого заряда на расстоянии $r$ рассчитывается по формуле $E=k\frac{Q}{r^2}=\frac1{4\pi\varepsilon_0}\cdot\frac Q{r^2}.$ Следовательно, на расстоянии $2R$ модуль напряженности $E= \frac{ 1}{ 4 \pi \varepsilon _{0} } \cdot \frac{ Q}{ 4R ^2 } .$\par
Б) Потенциал по сферически симметричного распределения заряда вне этого заряда на расстоянии $r$ рассчитывается по формуле $\varphi=k \frac{ Q}{ r } = \frac{ 1}{ 4 \pi \varepsilon _{0} } \cdot \frac{ Q}{ r } .$ Значит, на поверхности шара, то есть на расстоянии $R$, потенциал $\varphi= \frac{ 1}{ 4 \pi \varepsilon _{0} } \cdot \frac{ Q}{ R } .$\par
Ответ: $43$.\par
}

% Source sdamgia-16857: https://phys-ege.sdamgia.ru/problem?id=16857
\BankTask{sdamgia-16857}{
На уединенной неподвижной проводящей сфере радиусом $R$ находится положительный заряд $Q$. Сфера находится в вакууме. Напряженность электростатического поля сферы в точке $A$ равна $36 \,\text{В/м}$. Все расстояния указаны на рисунке.\par
$1$. Потенциал электростатического поля в точке $A$ выше, чем в точке $F$: $\varphi_{A} > \varphi_{F}.$\par
$2$. Потенциал электростатического поля в точках $B$ и $D$ одинаков: $\varphi_{B} = \varphi_{D}.$\par
$3$. Потенциал электростатического поля в точках $A$ и $B$ одинаков: $\varphi_{A} = \varphi_{B}.$\par
$4$. Напряженность электростатического поля в точке $C$ $E_{C}=9 \,\text{В/м}$.\par
$5$. Напряженность электростатического поля в точке $B$ $E_{B}=0$.\par

\begin{center}\begin{tikzpicture}[book diagram,>=stealth]
\draw(0,0)circle(1.4);\draw(-.7,0)--(2.8,0)(0,0)--(0,-2.8);\foreach\x/\y/\label in{1.4/0/D,2.8/0/A,0/-1.4/F,0/-2.8/C}{\fill(\x,\y)circle(1.5pt);\node[above right]at(\x,\y){$\label$};}\fill(-.7,0)circle(1.5pt)node[above left]{$B$};\fill(0,0)circle(1.5pt);\draw[<->](-.7,0)--(0,0)node[midway,above]{$R/2$};\draw[<->](0,0)--(1.4,0)node[midway,above]{$R$};\draw[<->](1.4,0)--(2.8,0)node[midway,above]{$R$};\draw[<->](0,0)--(0,-1.4)node[midway,left]{$R$};\draw[<->](0,-1.4)--(0,-2.8)node[midway,left]{$R$};
\end{tikzpicture}\end{center}
\par
}{
\smallskip\noindent\textbf{Решение}\par
Для точечного заряда и поля снаружи заряженной сферы верны формулы: $\varphi = k \frac{ Q}{ r } ,E=k \frac{ Q}{ r ^2 } ,$ где $Q$ — величина точечного заряда или заряда сферы, $r$ — расстояние до точечного заряда или центра сферы. Потенциал на границе сферы равен потенциалу любой точки внутри сферы. Напряженность поля внутри сферы равна нулю.\par
Рассмотрим предложенные утверждения.\par
$1$. Расстояние от точки $A$ до центра сферы больше, чем расстояние от центра сферы до точки $F$, следовательно, потенциал в точке $F$ больше потенциала в точке $A$. Утверждение $1$ неверно.\par
$2$. Как известно из сказанного выше, потенциал внутри сферы равен потенциалу на границе сферы. Утверждение $2$ верно.\par
$3$. Потенциал в точке $B$ равен потенциалу на поверхности сферы, расстояние от точки $A$, до центра сферы больше радиуса сферы, следовательно, потенциал в точке $A$ меньше потенциала в точке $B$. Утверждение $3$ неверно.\par
$4$. Заметим, что точки $A$ и $C$ находятся на одинаковом расстоянии от центра сферы, следовательно, напряженности поля в этих точках равны и равны $36 \,\text{В/м}$. Утверждение $4$ неверно.\par
$5$. Как известно из сказанного выше, напряженность поля внутри сферы равна нулю. Утверждение $5$ верно.\par
Ответ: $25$.\par
}

% Source sdamgia-31415: https://phys-ege.sdamgia.ru/problem?id=31415
\BankTask{sdamgia-31415}{
Стеклянные незаряженные кубики~$1$ и $2$ и алюминиевые незаряженные кубики~$3$ и $4$ сблизили вплотную и поместили в электрическое поле, напряженность которого направлена горизонтально влево, как показано в левой части рисунка. Затем кубики раздвинули и уже после этого выключили электрическое поле (правая часть рисунка).\par
$1$. В электрическом поле (левая часть рисунка) часть электронов перешла из кубика~$4$ в кубик~$3$.\par
$2$. Кубик $2$, изображенный в правой части рисунка, не заряжен.\par
$3$. В электрическом поле (левая часть рисунка) правая грань кубика~$2$ заряжена положительно.\par
$4$. Кубик $3$, изображенный в правой части рисунка, имеет положительный заряд.\par
$5$. В электрическом поле (левая часть рисунка) разность потенциалов между левой гранью кубика~$3$ и правой гранью кубика~$4$ равна нулю.\par

\begin{center}\begin{tikzpicture}[book diagram,>=stealth]
\draw(0,0)rectangle(1.2,1.2);\draw(.6,0)--(.6,1.2)(0,.6)--(1.2,.6);\foreach\x/\y/\s in{.3/.9/1,.9/.9/2,.3/.3/3,.9/.3/4}\node at(\x,\y){$\s$};\foreach\y in{.3,.5,.7,.9}\draw[->](-.2,\y)--(-1,\y);\node[above]at(-.7,1){$\vec E$};\draw[dashed](1.65,-.2)--(1.65,1.4)(1.8,-.2)--(1.8,1.4);\foreach\x/\y/\s in{2.1/.8/1,3.1/.8/2,2.1/-.2/3,3.1/-.2/4}{\draw(\x,\y)rectangle++(.6,.6);\node at(\x+.3,\y+.3){$\s$};}
\end{tikzpicture}\end{center}
\par
}{
\smallskip\noindent\textbf{Решение}\par
$1$. Неверно. Электроны под действием электрического поля движутся против силовых линий, поэтому переместятся вправо, тогда часть электронов перейдет с кубика~$3$ в кубик~$4$.\par
$2$. Верно. Стекло является диэлектриком, в котором нет свободных носителей заряда. Поэтому кубики из стекла останутся незаряженными.\par
$3$. Неверно. В кубиках из стекла произойдет поляризация, в результате которой правая часть кубика~$2$ будет иметь отрицательный заряд.\par
$4$. Верно. Поскольку электроны в металлах переместятся против силовых линий, то в кубике~$3$ будет недостаток электронов, он будет заряжен положительно.\par
$5$. Верно. В проводнике, помещенном в электростатическое поле, напряженность станет равной $0$, поэтому и разность потенциалов между левой гранью кубика~$3$ и правой гранью кубика~$4$ равна нулю.\par
Ответ: $245$.\par
}

% Source sdamgia-6067: https://phys-ege.sdamgia.ru/problem?id=6067
\BankTask{sdamgia-6067}{
Если потереть шерстью эбонитовую палочку, то она электризуется, приобретая отрицательный заряд, и стрелка электрометра при поднесении палки к его шару отклоняется, а при удалении палки — возвращается к неотклоненному состоянию. Если же в момент поднесения наэлектризованной палки к электрометру коснуться рукой его металлического корпуса и сразу же убрать руку, то после удаления палки отклонение стрелки сохраняется, хотя и меньшее по величине.\par
Объясните, основываясь на известных физических законах и закономерностях, почему это происходит.\par
Электрометр (см. рис.) представляет собой металлический цилиндрический корпус, передняя и задняя стенки которого стеклянные. Корпус закреплен на изолирующей подставке. Через изолирующую втулку внутрь корпуса сверху входит металлическая трубка, заканчивающаяся внизу стержнем с установленной на нем легкоподвижной стрелкой, отклонение которой определяется величиной заряда. Стрелка может вращаться вокруг горизонтальной оси. Внутри корпуса установлена шкала электрометра, по которой определяется отклонение стрелки. Снаружи корпуса, наверху трубки прикрепляется металлический шар или тарелка, к которой подносят заряженные тела.\par

\begin{center}\begin{tikzpicture}[book diagram,>=stealth]
\begin{scope}[shift={(0,0)}]\draw(0,0)circle(.65);\draw(0,.65)--(0,1.1);\draw(0,1.3)circle(.2);\draw(0,-.65)--(0,-1)--(-.4,-1)(0,-1)--(.4,-1);\draw(0,.35)--(0,-.35);\draw[thick](0.4,.48)--(-0.4,-.48);\foreach\a in{220,235,...,280}\draw({.48*cos(\a)},{.48*sin(\a)})--({.58*cos(\a)},{.58*sin(\a)});\end{scope}\draw[thick](.4,1.7)--(1.7,1.2);
\end{tikzpicture}\end{center}
\par
}{
\smallskip\noindent\textbf{Решение}\par
$1$. При поднесении наэлектризованной отрицательным зарядом эбонитовой палки к шару электрометра, в силу явления электростатической индукции и закона сохранения заряда в изолированной системе тел, шар и стрелка электрометра заряжаются разноименными и равными по величине зарядами (шар — «+», стрелка — «-»). При этом часть металлического корпуса электрометра вблизи шкалы заряжается положительным зарядом в силу того же явления электростатической индукции, а остальная часть — равным ему по величине в силу закона сохранения заряда в изолированной системе тел отрицательным зарядом.\par
$2$. Стрелка электрометра отклоняется, так как одноименные заряды на стрелке и на стержне отталкиваются, а разноименные заряды на стрелке и на корпусе электрометра притягиваются, согласно закону взаимодействия зарядов.\par
$3$. При удалении наэлектризованной палки от электрометра одинаковые индуцированные заряды разных знаков на его шаре и на стрелке, а также меньшие по величине заряды на корпусе электрометра вблизи его шкалы и вдали от стрелки компенсируются, и отклонение стрелки прекращается.\par
$4$. Если коснуться корпуса электрометра рукой после поднесения к нему наэлектризованной палки и сразу убрать руку, то часть индуцированных на корпусе зарядов (отрицательных) стечет на прикоснувшегося к нему\par
человека, и на корпусе электрометра останется нескомпенсированный положительный заряд.\par
$5$. В силу явления электростатической индукции после удаления палки этот положительный нескомпенсированный заряд на корпусе электрометра вызовет появление отрицательного заряда на стрелке электрометра, расположенной вблизи шкалы, и положительного — на шаре электрометра, что и приведет к отклонению стрелки, хотя и меньшему, чем при поднесении заряженной эбонитовой палки к электрометру.\par
}

% Source sdamgia-38948: https://phys-ege.sdamgia.ru/problem?id=38948
\BankTask{sdamgia-38948}{
На столе установили два незаряженных электрометра и соединили их металлическим стержнем с изолирующей ручкой (рис. $1$). Затем к первому электрометру поднесли, не касаясь его шара, положительно заряженную палочку (рис. $2$). Не убирая палочки, сняли стержень, а затем убрали палочку.\par
Ссылаясь на известные вам законы и явления, объясните, почему электрометры оказались заряженными, и определите знаки заряда каждого из электрометров после того, как палочку убрали.\par

\begin{center}\begin{tikzpicture}[book diagram,>=stealth]
\begin{scope}[shift={(0,0)}]\draw(0,0)circle(.65);\draw(0,.65)--(0,1.1);\draw(0,1.3)circle(.2);\draw(0,-.65)--(0,-1)--(-.4,-1)(0,-1)--(.4,-1);\draw(0,.35)--(0,-.35);\draw[thick](0,.48)--(0,-.48);\foreach\a in{220,235,...,280}\draw({.48*cos(\a)},{.48*sin(\a)})--({.58*cos(\a)},{.58*sin(\a)});\end{scope}\begin{scope}[shift={(2.4,0)}]\draw(0,0)circle(.65);\draw(0,.65)--(0,1.1);\draw(0,1.3)circle(.2);\draw(0,-.65)--(0,-1)--(-.4,-1)(0,-1)--(.4,-1);\draw(0,.35)--(0,-.35);\draw[thick](0,.48)--(0,-.48);\foreach\a in{220,235,...,280}\draw({.48*cos(\a)},{.48*sin(\a)})--({.58*cos(\a)},{.58*sin(\a)});\end{scope}\draw(0,1.5)--(2.4,1.5);\draw[fill=white](1.1,1.35)rectangle(1.25,1.85);\node[below]at(0,-1){$1$};\node[below]at(2.4,-1){$2$};\node[below]at(1.2,-1.5){Рис. 1};
\end{tikzpicture}\end{center}
\par

\begin{center}\begin{tikzpicture}[book diagram,>=stealth]
\begin{scope}[shift={(0,0)}]\draw(0,0)circle(.65);\draw(0,.65)--(0,1.1);\draw(0,1.3)circle(.2);\draw(0,-.65)--(0,-1)--(-.4,-1)(0,-1)--(.4,-1);\draw(0,.35)--(0,-.35);\draw[thick](.52,.3)--(-.52,-.3);\foreach\a in{220,235,...,280}\draw({.48*cos(\a)},{.48*sin(\a)})--({.58*cos(\a)},{.58*sin(\a)});\end{scope}\begin{scope}[shift={(2.4,0)}]\draw(0,0)circle(.65);\draw(0,.65)--(0,1.1);\draw(0,1.3)circle(.2);\draw(0,-.65)--(0,-1)--(-.4,-1)(0,-1)--(.4,-1);\draw(0,.35)--(0,-.35);\draw[thick](.52,.3)--(-.52,-.3);\foreach\a in{220,235,...,280}\draw({.48*cos(\a)},{.48*sin(\a)})--({.58*cos(\a)},{.58*sin(\a)});\end{scope}\draw(0,1.5)--(2.4,1.5);\draw[fill=white](1.1,1.35)rectangle(1.25,1.85);\draw[thick](-1.1,1.7)--(-.25,2.1);\node[above]at(-.7,2){$+$};\node[below]at(0,-1){$1$};\node[below]at(2.4,-1){$2$};\node[below]at(1.2,-1.5){Рис. 2};
\end{tikzpicture}\end{center}
\par
}{
\smallskip\noindent\textbf{Решение}\par
Вокруг положительно заряженной палочки существует электрическое поле. Под его действием свободные электроны в проводнике (металлической палочке, металлических шарах) придут в движение, причем сила электрического действия будет направлена к заряженной палочке. В результате на шаре электрометра $1$ возникнет избыток электронов, а на шаре электрометра $2$ возникнет недостаток электронов. Данный процесс будет происходить до тех пор, пока потенциалы всех точек не станут равными. Таким образом, электрометр $1$ будет заряжен отрицательно, а электрометр $2$ — положительно.\par
Электрометры образуют изолированную систему, в которой выполняется закон сохранения электрического заряда. Так как до поднесения палочки система не имела заряда, то после на каждом шаре заряды будут равными по модулю и противоположными по знаку. После удаления стержня и палочки на каждом шаре останется индуцированный заряд.\par
Ответ: электрометр $1$ имеет отрицательный заряд, электрометр $2$ — положительный.\par
}

\end{document}
